\section{Word2Vec 算法详解}

\subsection{Word2Vec 的核心思想}

Word2Vec 是 Tomas Mikolov 等人于 2013 年在 Google 提出的词向量学习方法。虽然不是第一个词向量算法（词向量的研究可追溯到 2000 年左右），但因其\textbf{简洁高效}而广受欢迎。

\begin{figure}[H]
\centering
\includegraphics[width=0.95\textwidth]{slides-images/slide-025.jpg}
\caption{Word2Vec 概览：给定中心词 ``into''，计算窗口内各上下文词出现的概率\protect\footnotemark}
\end{figure}
\footnotetext{Slides 第26页。}

Word2Vec（Skip-gram 模型）的基本思想：

\begin{importantbox}{Word2Vec Skip-gram 模型}
\begin{enumerate}
    \item 准备一个大型文本语料库（corpus）
    \item 为词表中每个词分配一个随机初始化的向量
    \item 遍历语料库中的每个位置 $t$，取当前位置的词为\textbf{中心词}（center word），取固定窗口 $m$ 内的词为\textbf{上下文词}（context/outside words）
    \item 利用中心词和上下文词的向量相似度，计算上下文词的出现概率
    \item 不断调整所有词向量，使得观察到的上下文的概率尽可能大
\end{enumerate}
\end{importantbox}

\begin{knowledgebox}{语料库（Corpus）}
Corpus 是拉丁语 ``body''（身体）的意思，在 NLP 中指一个大型文本集合。注意其复数形式是 \textbf{corpora}（不是 ``corpuses'' 或 ``corpi''）。Manning 在课上特别提醒：如果在作业中写了错误的复数形式，说明没有认真听第一节课！
\end{knowledgebox}

\subsection{似然函数与目标函数}

\subsubsection{似然函数}

对于语料库中的每个位置 $t$（$t = 1, \ldots, T$），给定中心词 $w_t$，我们希望在窗口大小为 $m$ 的范围内，上下文词 $w_{t+j}$（$-m \le j \le m$，$j \ne 0$）出现的概率尽可能高。整体似然函数为：

$$L(\theta) = \prod_{t=1}^{T} \prod_{\substack{-m \le j \le m \\ j \ne 0}} P(w_{t+j} \mid w_t; \theta)$$

\begin{itemize}
    \item $T$：语料库中词的总数
    \item $m$：上下文窗口大小（Manning 的示例中 $m=2$）
    \item $\theta$：模型的所有参数（即所有词向量）
\end{itemize}

\subsubsection{从似然到损失函数}

将似然函数转化为可优化的损失函数需要两步变换：

\begin{enumerate}
    \item \textbf{取负号}：将最大化问题转为最小化问题（对应梯度下降）
    \item \textbf{取对数}：将乘积转为求和，简化数学处理
\end{enumerate}

最终得到\textbf{平均负对数似然}（average negative log-likelihood）作为目标函数：

$$J(\theta) = -\frac{1}{T} \sum_{t=1}^{T} \sum_{\substack{-m \le j \le m \\ j \ne 0}} \log P(w_{t+j} \mid w_t; \theta)$$

\begin{figure}[H]
\centering
\includegraphics[width=0.95\textwidth]{slides-images/slide-028.jpg}
\caption{Word2Vec 目标函数：最小化平均负对数似然，每个词使用两个向量（中心向量 $v_w$ 和上下文向量 $u_w$）\protect\footnotemark}
\end{figure}
\footnotetext{Slides 第29页。}

\begin{knowledgebox}{为什么取对数？}
\begin{itemize}
    \item \textbf{数学便利}：$\log$ 将乘积变为求和，简化求导过程
    \item \textbf{数值稳定}：大量小概率相乘会导致数值下溢，取对数后变为负数的求和
    \item \textbf{单调性}：$\log$ 是单调递增函数，最大化 $\log L$ 等价于最大化 $L$
\end{itemize}
\end{knowledgebox}

\subsection{Softmax 概率计算}

\subsubsection{两套向量的设计}

Word2Vec 的一个关键设计是：每个词拥有\textbf{两个}向量：

\begin{itemize}
    \item $v_w$：当 $w$ 作为\textbf{中心词}（center word）时使用
    \item $u_w$：当 $w$ 作为\textbf{上下文词}（context/outside word）时使用
\end{itemize}

这一设计使得数学推导更简洁。

\subsubsection{Softmax 公式}

给定中心词 $c$ 和上下文词 $o$，条件概率定义为：

$$P(o \mid c) = \frac{\exp(u_o^T v_c)}{\sum_{w \in V} \exp(u_w^T v_c)}$$

\begin{itemize}
    \item $u_o^T v_c$：上下文词向量与中心词向量的\textbf{点积}，衡量二者的相似度
    \item $\exp(\cdot)$：指数函数，将任意实数映射为正数
    \item $\sum_{w \in V}$：对词表中所有词求和，起\textbf{归一化}作用，确保概率之和为 1
\end{itemize}

\begin{figure}[H]
\centering
\includegraphics[width=0.95\textwidth]{slides-images/slide-029.jpg}
\caption{Softmax 函数详解：先做点积，再取指数确保正数，最后归一化得到概率分布\protect\footnotemark}
\end{figure}
\footnotetext{Slides 第30页。}

\begin{importantbox}{Softmax 函数的直觉理解}
Softmax 是深度学习中将任意实数向量转化为概率分布的标准工具：
\begin{enumerate}
    \item \textbf{点积}$u_o^T v_c$：度量两个词向量的相似度（无界实数）
    \item \textbf{取指数}$\exp(\cdot)$：将实数映射为正数（确保概率非负）
    \item \textbf{归一化}（除以总和）：确保所有概率之和为 1
\end{enumerate}
名字中的 ``soft'' 表示它是 ``max'' 的平滑版本：它\textbf{放大}了最大项的概率，但仍然给较小项分配非零概率。
\end{importantbox}

\subsection{模型参数与优化}

\subsubsection{参数空间}

Word2Vec 的\textbf{全部}参数就是所有词的向量。假设词表大小为 $V$，词向量维度为 $d$，则：

\begin{figure}[H]
\centering
\includegraphics[width=0.95\textwidth]{slides-images/slide-030.jpg}
\caption{参数向量 $\theta$：由所有词的中心向量 $v_w$ 和上下文向量 $u_w$ 拼接而成，维度为 $\mathbb{R}^{2dV}$\protect\footnotemark}
\end{figure}
\footnotetext{Slides 第31页。}

$$\theta = \begin{bmatrix} v_{\text{aardvark}} \\ v_a \\ \vdots \\ v_{\text{zebra}} \\ u_{\text{aardvark}} \\ u_a \\ \vdots \\ u_{\text{zebra}} \end{bmatrix} \in \mathbb{R}^{2dV}$$

例如，当 $V = 400{,}000$、$d = 100$ 时，总参数量为 $2 \times 400{,}000 \times 100 = 80{,}000{,}000$（8千万）。

\subsubsection{梯度下降优化}

优化过程的核心是\textbf{梯度下降}（gradient descent）：

\begin{enumerate}
    \item 随机初始化所有词向量
    \item 计算目标函数 $J(\theta)$ 关于所有参数的梯度 $\nabla_\theta J(\theta)$
    \item 沿梯度的反方向更新参数：$\theta \leftarrow \theta - \alpha \nabla_\theta J(\theta)$
    \item 重复步骤 2--3 直到收敛
\end{enumerate}

这个过程类似于在高维空间中``走下坡路''：梯度指出了当前位置上升最快的方向，我们朝反方向走一小步，就能降低损失。

\subsection{梯度推导}

Manning 在课上详细推导了目标函数关于中心词向量 $v_c$ 的梯度。这是理解 Word2Vec 训练过程的关键。

\subsubsection{分解对数概率}

首先，我们需要计算单个 log 概率项的梯度：

$$\frac{\partial}{\partial v_c} \log P(o \mid c) = \frac{\partial}{\partial v_c} \log \frac{\exp(u_o^T v_c)}{\sum_{w \in V} \exp(u_w^T v_c)}$$

利用对数的性质 $\log \frac{a}{b} = \log a - \log b$，将其分解为两项：

$$= \underbrace{\frac{\partial}{\partial v_c} \log \exp(u_o^T v_c)}_{\text{项 1：分子}} - \underbrace{\frac{\partial}{\partial v_c} \log \sum_{w \in V} \exp(u_w^T v_c)}_{\text{项 2：分母}}$$

\begin{figure}[H]
\centering
\includegraphics[width=0.95\textwidth]{slides-images/slide-033.jpg}
\caption{梯度推导过程（板书）：将对数概率分解为分子和分母两项分别求导\protect\footnotemark}
\end{figure}
\footnotetext{Slides 第34页。}

\subsubsection{项 1 的求导}

分子部分非常简单，$\log$ 和 $\exp$ 互为逆运算，直接消掉：

$$\frac{\partial}{\partial v_c} \log \exp(u_o^T v_c) = \frac{\partial}{\partial v_c} (u_o^T v_c) = u_o$$

\begin{knowledgebox}{向量点积的导数}
对于向量点积 $u_o^T v_c = \sum_{i=1}^{d} (u_o)_i (v_c)_i$，关于 $v_c$ 的偏导数为：

$$\frac{\partial}{\partial (v_c)_j} \sum_{i=1}^{d} (u_o)_i (v_c)_i = (u_o)_j$$

因为只有 $i = j$ 的项对 $(v_c)_j$ 有贡献。将所有分量拼起来，得到向量 $u_o$。

这不是高中单变量微积分——这里的导数是一个\textbf{向量}！
\end{knowledgebox}

\subsubsection{项 2 的求导（链式法则）}

分母部分需要两次应用\textbf{链式法则}（chain rule）：

\textbf{第一次链式法则}（$\log$ 的导数 $\times$ 内部函数的导数）：

$$\frac{\partial}{\partial v_c} \log \sum_{w} \exp(u_w^T v_c) = \frac{1}{\sum_{w} \exp(u_w^T v_c)} \cdot \frac{\partial}{\partial v_c} \sum_{w} \exp(u_w^T v_c)$$

\textbf{第二次链式法则}（$\exp$ 的导数 $\times$ 内部函数的导数）：

$$\frac{\partial}{\partial v_c} \sum_{w} \exp(u_w^T v_c) = \sum_{w} \exp(u_w^T v_c) \cdot u_w$$

\subsubsection{合并得到最终结果}

将两项合并：

$$\frac{\partial}{\partial v_c} \log P(o \mid c) = u_o - \frac{\sum_{x=1}^{V} \exp(u_x^T v_c) \cdot u_x}{\sum_{w=1}^{V} \exp(u_w^T v_c)}$$

\begin{figure}[H]
\centering
\includegraphics[width=0.95\textwidth]{slides-images/slide-035.jpg}
\caption{梯度推导最终结果（板书）：``观察值 $-$ 期望值''的形式\protect\footnotemark}
\end{figure}
\footnotetext{Slides 第36页。}

注意到分数部分恰好是 softmax 概率 $P(x \mid c)$，因此最终结果可以优雅地写为：

$$\boxed{\frac{\partial}{\partial v_c} \log P(o \mid c) = u_o - \sum_{x=1}^{V} P(x \mid c) \cdot u_x}$$

\begin{importantbox}{``观察值 $-$ 期望值'' 的直觉}
梯度具有非常直观的含义：
\begin{itemize}
    \item $u_o$：\textbf{观察到的}上下文词向量（实际出现的词）
    \item $\sum_{x} P(x \mid c) \cdot u_x$：\textbf{期望的}上下文词向量（模型预测的加权平均）
\end{itemize}
梯度 = 观察值 $-$ 期望值。这意味着：
\begin{itemize}
    \item 如果模型的预测完美匹配观察数据，梯度为零——已达到最优
    \item 否则，梯度会驱动参数调整，使模型的预测更接近实际观察
\end{itemize}
这种``observed $-$ expected''的形式在统计学和机器学习中非常常见。
\end{importantbox}

\begin{warningbox}{计算复杂度问题}
在上述公式中，归一化项 $\sum_{w \in V} \exp(u_w^T v_c)$ 需要遍历整个词表（可能有几十万个词）。这使得每次梯度计算的复杂度为 $O(V)$，在实际中非常昂贵。后续课程将介绍 \textbf{负采样}（negative sampling）等技巧来解决这个问题。
\end{warningbox}

\subsection{随机梯度下降（SGD）}

由于目标函数 $J(\theta)$ 是对语料库中\textbf{所有}窗口的求和（可能有数十亿个），计算完整梯度 $\nabla_\theta J(\theta)$ 的代价极其昂贵。

\begin{figure}[H]
\centering
\includegraphics[width=0.95\textwidth]{slides-images/slide-038.jpg}
\caption{随机梯度下降（SGD）：每次只采样一个窗口就更新参数，大幅提升效率\protect\footnotemark}
\end{figure}
\footnotetext{Slides 第39页。}

解决方案是\textbf{随机梯度下降}（Stochastic Gradient Descent, SGD）：

\begin{lstlisting}[caption={SGD / Mini-batch Gradient Descent 伪代码},language=Python]
while True:
    window = sample_window(corpus)
    theta_grad = evaluate_gradient(J, window, theta)
    theta = theta - alpha * theta_grad
\end{lstlisting}

\begin{importantbox}{SGD 的核心思想}
\begin{itemize}
    \item 不计算完整梯度，而是每次\textbf{随机采样}一个（或一小批）窗口
    \item 基于这个小样本估计梯度并更新参数
    \item 虽然每次更新的方向有噪声，但\textbf{平均而言}会朝正确方向前进
    \item 实践中通常使用 \textbf{mini-batch}（小批量）而非单个样本，以兼顾效率和稳定性
\end{itemize}
\end{importantbox}

\subsection{本章小结}

Word2Vec 的训练流程可以概括为：

\begin{enumerate}
    \item \textbf{初始化}：为每个词随机生成中心向量和上下文向量
    \item \textbf{遍历语料}：对每个位置，取中心词和窗口内的上下文词
    \item \textbf{计算概率}：通过 softmax 计算上下文词的条件概率
    \item \textbf{计算梯度}：利用``观察值 $-$ 期望值''公式
    \item \textbf{更新参数}：通过 SGD 调整词向量
    \item \textbf{重复}直到收敛
\end{enumerate}

最终得到的词向量能够编码丰富的语义信息，使得语义相似的词在向量空间中彼此接近。

%% ========== 总结与延伸 ==========

